  We propose a unified structural framework for superconductivity
  that applies across conventional and strongly correlated regimes.
  In this approach, pairing corresponds to the stabilization and
  projective phase locking of a topological composite class ($w=2$)
  that minimizes a real-space frustration functional under lattice constraints.
  Rather than taking a material-specific pairing interaction as the primary
  organizing principle, the mechanism identifies superconductivity with the
  emergence of a stable composite sector whose collective phase coherence
  generates the London and Ginzburg--Landau structure.

  In conventional superconductors, this composite class is stabilized
  by phonon-mediated attraction, consistent with Bardeen--Cooper--Schrieffer theory.
  In strongly correlated systems, stabilization instead arises from
  frustration reduction in real space, naturally separating amplitude
  and phase scales and providing a geometric account of pseudogap phenomenology.

  The symmetry of the order parameter is selected by a real-space
  minimization of the fiber raccordement cost.
  The framework predicts $d_{x^2-y^2}$ symmetry in strongly staggered
  cuprates and extended $s$-wave symmetry in nickelates,
  where antiferromagnetic frustration is reduced and partially isotropized.
  To leading order, the critical temperature scales as
  $T_c \sim \delta J$, linking phase stiffness directly
  to independently measurable exchange and frustration amplitudes.

  The theory yields falsifiable predictions for gap symmetry,
  phase stiffness scaling, disorder sensitivity,
  and pressure dependence in nickelates.
  It provides a constraint-based geometric description
  of low- and high-temperature superconductivity
  within a common structural mechanism.
